\part{Факторкольца и иже с ними}

\section{Кольца}

\subsection{Аксиомы поля}

\begin{definition}
	\tdef{Поле}~"--- это тройка $(F, +, \cdot)$, состоящая из множества~$F$ и двух бинарных операций
	\begin{align*}
		+ \colon
		F \times F &\to F
		,
		\\
		\cdot \colon
		F \times F &\to F
		,
	\end{align*}
	удовлетворяющих следующим свойствам.
	\refstepcounter{equation}%
	\[
		\begin{array}{r@{}r@{{}={}}lll}
			\multicolumn{5}{c}{%
				\text{Свойства сложения:}
			}
			\\
			\forall a,\ b,\ c \in F
			\holds {}
			&
			(a + b) + c & a + (b + c)
			&
			\text{(ассоциативность);}
			&
			(\arabic{equation}.1.1)%
			% \label{eq:fld:p:ass}
			\\
			\exists 0 \in F
			\suchthat {}
			&
			a + 0 & a
			&
			\text{(нейтральный);}
			&
			(\arabic{equation}.1.2)%
			% \label{eq:fld:p:ntr}
			\\
			\forall a \in F\ 
			\exists \ta \in F\ 
			\suchthat {}
			&
			a + \ta & 0
			&
			\text{(обратный);}
			&
			(\arabic{equation}.1.3)%
			% \label{eq:fld:p:inv}
			\\
			\forall a,\ b \in F
			\holds {}
			&
			a + b & b + a
			&
			\text{(коммутативность);}
			&
			(\arabic{equation}.1.4)%
			% \label{eq:fld:p:cmt}
			\\\\
			\multicolumn{5}{c}{%
				\text{Свойства умножения:}
			}
			\\
			\forall a,\ b,\ c \in F
			\holds {}
			&
			(a \cdot b) \cdot c & a \cdot (b \cdot c)
			&
			\text{(ассоциативность);}
			&
			(\arabic{equation}.2.1)%
			% \label{eq:fld:m:ass}
			\\
			\exists 1 \in F
			\suchthat {}
			&
			a \cdot 1 & a
			&
			\text{(нейтральный);}
			&
			(\arabic{equation}.2.2)%
			% \label{eq:fld:m:ntr}
			\\
			\forall a \in F^\times \ 
			\exists \ha \in F^\times \ 
			\suchthat {}
			&
			a \cdot \ha & 1
			&
			\text{(обратный);}
			&
			(\arabic{equation}.2.3)
			% \label{eq:fld:m:inv}
			\\
			\forall a,\ b \in F
			\holds {}
			&
			a \cdot b & b \cdot a
			&
			\text{(коммутативность);}
			&
			(\arabic{equation}.2.4)
			% \label{eq:fld:m:cmt}
			\\\\
			\multicolumn{5}{c}{%
				\text{Связь сложения и умножения:}
			}
			\\
			\forall a,\ b,\ c \in F
			\holds {}
			&
			a \cdot (b + c) &
			a \cdot b + a \cdot c
			&
			\text{(дистрибутивность);}
			&
			(\arabic{equation}.3.1)%
			% \label{eq:fld:distr}
			\\
			&
			\multicolumn{2}{c}{\mkern -7.9mu 0 \neq 1}
			&
			\text{(невырожденность).}
			&
			(\arabic{equation}.3.2)%
			% \label{eq:fld:ndeg}
			% \label{<+label+>}
		\end{array}
	\]
\end{definition}

Свойства
%\ref{eq:fld:p:ass}--\ref{eq:fld:p:cmt}
(\arabic{equation}.1.1)--(\arabic{equation}.1.4)~"--- это определение группы $(F, +)$ по сложению.

Свойства (\arabic{equation}.2.1)--(\arabic{equation}.2.4)~"--- это \emph{почти} определение группы $(F, \cdot)$ по умножению, за исключением того, что обратный по умножению определён только для ненулевого элемента.

Множество $F^\times := F \setminus \set*{0}$ называется \emph{мультипликативной подгруппой} поля~$F$, так как если сузить операции $+$ и $\cdot$ на $F^\times \times F^\times$, то свойства~(\arabic{equation}.2.1)--(\arabic{equation}.2.4) дадут определение группы~$(F^\times, \cdot)$.

Если $0 = 1$ (в смысле равенства нейтральных элементов по сложению и умножению), то множество~$F$ тоже называется полем~"--- но это будет поле из одного элемента.

\begin{remark}
	Обычно элемент, обратный по сложению к~$a \in K$, обозначается~$-a$.

	А элемент, обратный по умножению к~$a$, обозначается~$a^{-1}$.
\end{remark}

\subsection{Аксиомы кольца}

\begin{definition}
	Уберём из аксиом поля все свойства умножения, кроме ассоциативности. Получим аксиомы \emph{(ассоциативного) кольца}%
\footnote{%
	В некоторых источниках от умножения не требуется даже ассоциативности~"--- и тогда кольцо с ассоциативным умножением называется \emph{ассоциативным кольцом}.%
}%
.
\end{definition}

\begin{definition}
	Если вдобавок к~(1.2.1) выполняется требование~(1.2.2) наличия нейтрального элемента, то кольцо называется \emph{(ассоциативным) кольцом с единицей}.
\end{definition}

\begin{definition}
	Если вдобавок к~(1.2.1) выполняется требование~(1.2.4) коммутативности по умножению, то кольцо называется \emph{(ассоциативным) коммутативным кольцом}.
\end{definition}

\begin{remark}
	По сложению \emph{любое} кольцо и так коммутативно, поэтому эпитет <<коммутативный>> при разговоре о кольцах относится только к умножению.
\end{remark}


\begin{definition}
	Некоммутативное поле называется \emph{телом}.
\end{definition}

\begin{remark}
	Даже если кольцо~"--- это коммутативное кольцо с единицей, в нём может быть не выполнено свойство~(1.2.3) обратимости \emph{любого} ненулевого элемента. При этом \emph{какие-то} элементы этого кольца могут быть обратимы.
\end{remark}

В кольцах возможны <<экзотические>> ситуации.
\begin{definition}
	Ненулевой элемент~$a \in K^\times$ кольца~$K$ называется \tdef{делителем нуля}, если существует ненулевой элемент~$b \in K^\times$, такой, что
	\[
		ab = 0
		.
	\]
\end{definition}

\begin{definition}
	Ненулевой элемент~$a$ кольца~$K$ называется \tdef{нильпотентом}, если $a^n = 0$ для некоторого~$n \in \NN$.
\end{definition}

\begin{definition}
	Кольцо с единицей и без делителей нуля называется \tdef{целостным}.
\end{definition}

\begin{definition}
	Кольцо с единицей и без нильпотентов называется \tdef{приведённым}.
\end{definition}


% \subsection{Определения}
% \begin{definition}
% 	\tdef{Кольцо}~"--- это тройка $(K, +, \cdot)$, включающая в себя множество~$K$ с двумя операциями
% 	\begin{align*}
% 		+ \colon
% 		K \times K &\to K
% 		,
% 		\\
% 		\cdot \colon
% 		K \times K &\to K
% 		,
% 	\end{align*}
% 	причём $(K, +)$~"--- группа по сложению.
% 	\label{<+label+>}
% \end{definition}

\subsection{Гомоморфизмы колец}

\begin{definition}
	Гомоморфизм колец $\phi \colon A \to B$~"--- это отображение подлежащих множеств, которое переводит сложение в сложение, а умножение на число~"--- в умножение на число: для любых $a'$, $a'' \in A$ выполняется
	\begin{align*}
		\phi(a' \plusA a'') &=
		\phi(a') \plusB \phi(a''),
		\\
		\phi(a' \cdotA a'') &=
		\phi(a') \cdotB \phi(a'').
	\end{align*}
\end{definition}

\newsavebox{\tempbox}
\savebox{\tempbox}{%
	\begin{minipage}{2in}
		\[
			\left(
				\begin{gathered}
					\phi(a') \plusB \phi(a'')
					\\
					\phi(a' \plusA a'')
				\end{gathered}
			\right)
		\]
	\end{minipage}
}
Это свойство можно выразить в виде коммутативной диаграммы
\[
	\xymatrix{
		A \times A
		\ar [r] _-{\plusA}
		\ar [d] _-{\phi \times \phi}
		&
		A
		\ar [d] ^-{\phi}
		\\
		B \times B
		\ar [r] _-{\plusB}
		&
		B
	}
	\qquad
	\xymatrix{
		(a', a'')
		\ar @{|->} [r] _-{\plusA}
		\ar @{|->} [d] _-{\phi \times \phi}
		&
		a' \plusA a''
		\ar @{|->} [d] ^-{\phi}
		\\
		\bigl( \phi(a'), \phi(a'') \bigr)
		\ar @{|->} [r] _-{\plusB}
		&
		\usebox{\tempbox}
	}
\]
(коммутативность означает, что можно двигаться по любому пути, указанном стрелками, и получить один и тот же результат).

Для операций умножения на число коммутативная диаграмма выглядит аналогично:
% \[
% 	\xymatrix{
% 		A \times A
% 		\ar [r] ^-{\cdotA}
% 		\ar [d] _-{\phi \times \phi}
% 		&
% 		A
% 		\ar [d] _-{\phi}
% 		\\
% 		B \times B
% 		\ar [r] ^-{\cdotB}
% 		&
% 		B
% 	}
% 	\qquad
% 	\xymatrix{
% 		(a', a'')
% 		\ar @{|->} [r] ^-{\cdotA}
% 		\ar @{|->} [d] _-{\phi \times \phi}
% 		&
% 		a' \cdotA a''
% 		\ar @{|->} [d] _-{\phi}
% 		\\
% 		\bigl( \phi(a'), \phi(a'') \bigr)
% 		\ar @{|->} [r] ^-{\cdotB}
% 		&
% 		\begin{aligned}
% 			\phi(a') &\cdotB \phi(a'')
% 			\\
% 			\phi(a' &\cdotA a'')
% 		\end{aligned}
% 	}
% \]

\begin{definition}
	Пусть $A$~"--- кольцо. Подмножество $B \subset A$ называется \emph{подкольцом}, если оно замкнуто относительно операций $\bound{+}_B$ и $\bound{\cdot}_B$~"--- ограничений соответстующих операций в~$A$. Иными словами, $B$~"--- подкольцо, если для всех $b_1$, $b_2 \in B$ выполняется
	\begin{align*}
		b_1 \cdot b_2 &\in B,
		\\
		b_1 + b_2 &\in B,
	\end{align*}
	где $+$ и $\cdot$~"--- операции в~$A$.
\end{definition}

Тот факт, что $B$~"--- подкольцо~$A$, записывается в виде~$B \leq A$.

\begin{definition}
	Пусть $K$ и $L$~"--- два подкольца в~$A$, $k \in K$. Будем писать
	\begin{align*}
		kL &:= \set{k\ell}{\ell \in L}
		=
		\bigcup_{\ell \in L} \set*{k\ell}
		,
		\\
		KL &:= \set{kL}{k \in K}
		=
		\bigcup_{k \in K} kL
		=
		\bigcup_{%
			\substack{
				k \in K
				\\
				\ell \in L
			}
		}
		\set*{k\ell}
		.
	\end{align*}
\end{definition}

\subsection{Факторкольца}

\begin{definition}
	Пусть $A$~"--- кольцо, а $I \leq A$~"--- его подкольцо. Это подкольцо называется \emph{идеалом}, если для любого~$a \in A$ выполняется
	\[
		aI \subset I
		.
	\]
\end{definition}

\begin{proposition}
	Если~$I \leq A$~"--- идеал в кольце~$A$, то элементы~$I$~"--- это всевозможные конечные суммы вида
	\[
		a_1 \theta_1 + \dots + a_n \theta_n,
	\]
	где $a_\nu \in A$, $\theta_\nu \in I$, $n \in \NN$.
\end{proposition}
\begin{proof}
	Действительно, если~$w \in I$, то для любого~$a \in A$
	\[
		aw \in I
	\]
	по определению идеала,
	а если~$w_1$ и~$w_2$~"--- два элемента идеала~$I$, то
	\[
		w_1 + w_2 \in I,
	\]
	так как~$I$~"--- подкольцо в~$A$.
\end{proof}

\begin{definition}
	Пусть $a$, $b \in \ZZ$; обозначим~$(a, b)$ идеал, порождённый~$a$ и~$b$, то есть
	\[
		(a, b) :=
		\set{ka + \ell b}{k, \ell \in \ZZ}
		.
	\]
	Наименьшее положительное (ненулевое) число~$d \in (a, b)$ называется \tdef{наибольшим общим делителем} чисел~$a$ и~$b$ и обозначается~$\gcd{a, b}$.
\end{definition}

% \begin{lemma}
% 	Идеал~$(a, b)$ состоит из чисел, кратных~$d$, и каждый общий делитель~$\delta$ элементов~$a$ и~$b$ входит в~$(a, b)$, то есть кратен~$d$.
% \end{lemma}
% \begin{proof}
% 	
% \end{proof}<++>



\begin{definition}
	Числа~$a$ и~$b \in \ZZ$ называются \tdef{взаимно простыми}, если $\gcd{a, b} = 1$. Обозначение: $a \mpr b$.
\end{definition}

По идеалам можно факторизовать.

\begin{definition}
	Пусть $A$~"--- кольцо, а~$I \leq A$~"--- идеал в нём. Тогда кольцо
	\[
		A/I :=
		\set{[a]}{a \in A},
	\]
	где
	\[
		[a] := a + I =
		\set{a + i}{i \in I} =
		\set{a' \in A}{a - a' \in I}
	\]
	с определёнными ниже операциями называется \emph{факторкольцом}.

	Множество~$[a]$ называется \emph{классом эквивалентности}, а $a$~"--- \emph{представителем} класса~$[a]$.
\end{definition}

\begin{remark}
	Подлежащее множество факторкольца~$A/I$ можно определить как фактормножество~$A/\sim$ по отношению эквивалентности
	\[
		(a' \sim a'') :=
		(a' - a'' \in I)
		.
	\]
\end{remark}

Опеределим операции сложения и умножения на элемент~$k \in A$ в факторкольце~$A/I$ (то есть превратим множество~$A/I$  в кольцо).

\begin{definition}
	Сложение в факторкольце~$A/I$ определим так:
	\[
		+ \colon
		\frac{A}{I} \times \frac{A}{I} \to \frac{A}{I}
		\colon
		([a], [b]) \mapsto [a + b]
		.
	\]
	Здесь сложение, записанное внутри квадратных скобок,~"--- это сложение в~$A$. Иными словами,
	\[
		[a] + [b] := [a + b].
	\]
	Аналогично определяем умножение на элемент~$A/I$:
	\[
		\cdot \colon
		\frac{A}{I} \times \frac{A}{I} \to \frac{A}{I}
		\colon
		([a], [b]) \mapsto [a \cdot b]
		.
	\]
	Иными словами,
	\[
		[a] \cdot [b] := [kb]
		.
	\]
\end{definition}

\begin{remark}
	Определение сложения и умножения в~$A/I$ \emph{корректно}, то есть отображения
	\begin{align*}
		+ \colon
		\frac{A}{I}
		\times
		\frac{A}{I}
		&\to
		\frac{A}{I}
		\\
		\cdot \colon
		\frac{A}{I}
		\times
		\frac{A}{I}
		&\to
		\frac{A}{I}
	\end{align*}
	--- действительно \emph{отображения}, то есть всюду определённые и функиональные отношения%
	\footnote{%
		Функциональность отношения~$\xi \subset X \times Y$ означает, что оно не расклеивает точки:
		\[
			\bigl( x_1 = x_2 \bigr)
			\implies
			\bigl(
				\xi(x_1) = \xi(x_2)
			\bigr)
		\]
		--- или, что то же самое,
		\[
			\bigl(
				\xi(x_1) \neq \xi(x_2)
			\bigr)
			\implies
			\bigl( x_1 \neq x_2 \bigr)
		\]
		(прообразы различных точек различны).
	}%
	.
	% \label{<+label+>}
\end{remark}
\begin{proofm}
	Покажем, что результат вычисления~$[a] \cdot [b]$ не зависит от выбора представителей классов~$[a]$ и~$[b]$:
	\begin{align*}
		[a] \cdot [b] &=
		(a + I) \cdot (b + I)
		=
		ab + aI + Ib + II
		\nlinea{=}
		ab + I + I + I
		=
		ab + I
		=:
		[ab]
		.
	\end{align*}
	Аналогично~"--- для сложения:
	\begin{align*}
		[a] + [b] &=
		(a + I) + (b + I)
		=
		a + b + I + I
		\nlinea{=}
		a + b + I
		=:
		[a+b]
		.
		\proofmark
	\end{align*}
\end{proofm}


\begin{definition}
	Гомоморфизм $\phi \colon A \to B$ называется \emph{изомофризмом}, если существует обратный гомоморфизм $\psi \colon B \to A$, такой, что $\phi \circ \psi = \id_{B}$, $\psi \circ \phi = \id_{A}$, где $\id_X$~"--- тождественный гомоморфизм на~$X$, переводящий~$x \in X$ в~$x$. Иными словами, для любых $a \in A$ и $b \in B$
	\begin{align*}
		(\psi \circ \phi)(a) &= a,
		\\
		(\phi \circ \psi)(b) &= b.
	\end{align*}
\end{definition}

\subsection{Кольца вычетов}

Пусть $n \in \NN$, $(n) := \set{nk}{k \in \ZZ}$~"--- идеал в~$\ZZ$. (Почему это идеал?) Определим
\begin{align*}
	\ZZ/(n) &:=
	\set{[a]}{a \in \ZZ},
	\\
	[a] &:= a + (n) =
	\set{a + kn}{k \in \ZZ}.
\end{align*}
Такая конструкция называется~\tdef{кольцом вычетов по модулю~$n$}.
При этом
\begin{align*}
	[a] + [b] &:= [a + b],
	\\
	[a] \cdot [b] &:=
	[a \cdot b].
\end{align*}

\begin{lemma}
	Так как для любого~$a \in \ZZ$ существует единственное разложение
	\[
		a = kn + r,
		\rlap{%
			\quad
			где
			\quad
			$
				\begin{system}
					&k \in \ZZ,
					\\
					&r \in [0, n) \cap \NN,
				\end{system}
			$
		}%
	\]
	то~$r$~"--- остаток от деления~$a$ на~$n$~"--- это наименьший положительный представитель класса~$[a] \in \ZZ/(n)$.
	% \label{<+label+>}
\end{lemma}

\begin{proof}
	Действительно,
	\[
		[a] = [kn + r] = [r],
	\]
	и если есть другой положительный представитель~$[r'] = [r]$, для которого $0 \leq r' < r$, то $r = r' + \ell n$. Но так как~$r = a - kn$, то $a - kn = r' + \ell n$, откуда $a - (k - \ell)n = r'$, причём $0 \leq r' < n$, и тогда~$r'$~"--- остаток. Противоречие.
\end{proof}

\begin{proposition}
	Если $n$~"--- не простое число, то есть $n = pq$ для некоторых~$p \in \NN\setminus\set*{0, 1}$ и~$q \in \NN \setminus \set*{0, 1}$, то в~$\ZZ/(n)$ есть делители нуля: 
	\[
		[p][q] = [pq] = [n] = [0].
	\]
\end{proposition}
\begin{example}
	В $\ZZ/(12)$ (циферблат часов) имеем $[3] \cdot [4] = [12] = [0]$.
\end{example}


\begin{proposition}
	Обратимость элемента~$[a] \in \ZZ/(n)$ равносильна тому, что $\gcd{a, n} = 1$.
\end{proposition}
\begin{proofm}
	Если существует~$[b] \in \ZZ/(n)$, такой, что $[a][b] = [1]$, то
	\begin{align*}
		[ab] &= 1,
		\\
		ab + (n) &= 1 + (n),
		\\
		ab + kn &= 1 + \ell n,
		\\
		ab + (k - \ell)n &= 1,
		\\
		ab + mn &= 1,
	\end{align*}
	а это равносильно утверждению~$\gcd{a, n} = 1$.
\end{proofm}

\begin{remark}
	Проверить, обратим ли данный класс~$[m] \in \ZZ/(n)$, и если да, то вычислить~$[m]^{-1}$, можно при помощи алгоритма Евклида.
\end{remark}

\begin{remark}
	Если~$p$~"--- простое, то для любого целого~$0 < a < p$ имеем
	\[
		\gcd{a, p} = 1
	,
	\]
	и поэтому любой ненулевой элемент обратим, а это значит, что~$\ZZ/(p)$~"--- поле.
\end{remark}

а

\begin{example}
	Вычислим остаток от деления~$2^{100}$ на~$5$. Этот остаток~"--- это наименьший положительный представитель класса~$\bigl[ 2^{100} \bigr] \in \ZZ/(5)$. Посчитать какого-нибудь представителя этого класса просто:
	\[
		\bigl[ 2^{100} \bigr] =
		\bigl[ 2^{2 \cdot 50} \bigr] =
		\bigl[ \bigl( 2^2 \bigr)^{50} \bigr] =
		[4]^{50}.
	\]
	А так как
	\[
		[4] = [4 + 0] = [4] + [0] = [4] + [-5] = [4 - 5] = [-1],
	\]
	то
	\[
		[4]^{50} = [-1]^{50} =
		\bigl[ (-1)^{50} \bigr] =
		[1]
		.
	\]
	Итак, $2^{100} \bmod 5 = 1$.
\end{example}

\subsection{Ядра гомоморфизмов}

% Пусть $A$ и $B$~"--- кольца с единицами $1_A \in A$ и $1_B \in B$.
Пусть $\phi \colon A \to B$~"--- гомоморфизм колец. Тогда \emph{ядро гомомофризма~$\phi$}~"--- это
\[
	\ker{\phi} := \phi^{-1}(0_B) =
	\set{a \in A}{\phi(a) = 0_B};
\]
\emph{образ гомоморфизма~$\phi$}~"--- это
\[
	\im{\phi} := \phi(A) = \set{\phi(a)}{a \in A} \subset B
	.
\]

\begin{proposition}
	Для любого гомоморфизма колец $\phi \colon A \to B$ ядро $\ker{\phi}$~"--- это подкольцо в~$A$, а образ $\im{\phi}$~"--- это подкольцо в~$B$.
	\label{<+label+>}
\end{proposition}
\begin{proofm}
	Покажем, что $\ker{\phi}$~"--- подкольцо в~$A$. Для этого покажем, что $\forall x,\ y \in \ker{\phi}$ выполняются условия $x \cdot y \in \ker{\phi}$.
	
	Пусть $x,\ y \in \ker{\phi}$. Это значит, что $\phi(x) = 0_B$, $\phi(y) = 0_B$. Тогда
	\[
		\phi(x \cdot y) :=
		\phi(x) \cdot \phi(y)
		=
		0_B \cdot 0_B = 0_B,
	\]
	и $x \cdot y \in \ker{\phi}$.

	\medskip

	Аналогично покажем, что~$\im{\phi} \leq B$. Пусть $p,\ q \in \im{B}$. Покажем, что тогда $pq \in \im{B}$.

	Так как $p\in \im{B}$, то существует (возможно, не единственный) элемент $x \in A$, такой, что $\phi(x) = p$; аналогично, для~$q \in \im{B}$ существует элемент $y \in A$, такой, что $\phi(y) = q$. Тогда элемент~$xy$~"--- это прообраз элемента~$pq$, так что $pq \in \im{B}$:
	\[
		\phi(xy) = \phi(x) \phi(y) = pq
		.
		\proofmark
	\]
\end{proofm}

\begin{theorem}[о гомоморфизме]
	Она повторяет теорему о гомоморфизме групп%
	\footnote{%
		Её иллюстрирует стишок:
		\begin{verse}
			Гомоморфный образ группы\\*
			До победы коммунизма\\*
			Изомофрен факторгруппе\\*
			По ядру гомоморфизма.
		\end{verse}
		Если $\psi \colon G \to H$~"--- гомоморфизм групп, то
		$
			\im{\psi} \cong
			G/\ker{\psi}
		$%
		.%
	}:
	если $\phi \colon A \to B$~"--- гомоморфизм колец, то
	\[
		\im{\phi} \cong
		\frac{A}{\ker{\phi}}
		.
	\]
\end{theorem}

\begin{lemma}
	Гомоморфизм полей~$\phi \colon F \to H$ переводит единицу в единицу.
\end{lemma}
\begin{proofm}
	Так как~$\phi$ уважает умножение,
	\[
		\phi(1_F) =
		\phi(1_F \cdot 1_F) =
		\phi(1_F) \cdot \phi(1_F) =
		\phi(1_F) \cdot \phi(1_F)
		.
	\]
	Так как~$H$ --- поле, то равенство
	\[
		\phi(1_F) = \phi(1_F) \cdot \phi(1_F)
	\]
	можно умножить на~$\bigl( \phi(1_F) \bigr)^{-1}$ и получить
	\[
		1_H = \phi(1_F).
		\proofmark
	\]
\end{proofm}

\begin{remark}
	Гомоморфизм колец с единицей~$\phi \colon K \to L$ не обязан переводить единицу в единицу.
\end{remark}
\begin{proof}
	Отображение
	\[
		\psi \colon
		K \to L
		\colon
		a \mapsto 0_L
	\]
	является гомоморфизмом колец (почему?), но отправляет~$1_K$ в~$0_L$.
\end{proof}

\subsection{Характеристика поля}

\begin{definition}
	Пусть $F$~"--- поле.

	Рассмотрим гомоморфизм колец $\iota \colon \ZZ \to F$. Это естественное вложение. Естественность означает, что $\iota(1) = 1_F$, где $1_F$~"--- единица поля~$F$. Так как~$\iota$~"--- гомоморфизм колец, то для любого целого~$n$
	\begin{align*}
		\iota(n)
		&=
		\iota(
			\underbrace{
				1 + 1 + \dots + 1
			}_{\text{$n$ раз}}
		)
		\nlinea{=}
		\underbrace{
			\iota(1) +
			\iota(1) +
			\dots +
			\iota(1)
		}_{\text{$n$ раз}}
		\nlinea{=}
		\underbrace{
			1_F +
			1_F +
			\dots +
			1_F
		}_{\text{$n$ раз}}
		.
	\end{align*}
	\tdef{Характеристика}~$\opchar{F}$ поля~$F$~"--- это наименьшее~$n \in \NN$, такое, что $\iota(n) = 0_F$.

	Если для любого~$n \in \NN$ имеем~$\iota(n) \neq 0$, то считаем~$\opchar{F} = 0$.
	\label{<+label+>}
\end{definition}

\begin{example}
	В поле~$F = \ZZ/(5)$ имеем~$\opchar{F} = 5$, так как
	\begin{align*}
		\iota(1) &= [1] \neq [0], \\
		\iota(2) &= [1] + [1] = [2] \neq [0], \\
		\iota(3) &= [1] + [1] + [1] = [3] \neq [0], \\
		\iota(4) &= [1] + [1] + [1] + [1] = [4] \neq [0], \\
		\iota(5) &= [1] + [1] + [1] + [1] + [1] = [5] = [0].
	\end{align*}

	Вообще для любого простого~$p \in \NN$ кольцо~$\FFp := \ZZ/(p)$ является полем, и~%
	$
		\opchar{\FFp}
		=
		p
	$.
	% \[
	% 	\opchar
	% 	\left( 
	% 		\frac{\ZZ}{(p)}
	% 	\right)
	% 	=
	% 	p.
	% \]

	Ещё более общий случай: для конечного поля~$\FF_{p^n}$ из~$p^n$ элементов, где~$p \in \NN$~"--- простое, а~$n \in \NN \setminus \set*{0}$, 
	$
		\opchar{\FF_{p^n}} = p^n
	$.
\end{example}

\section{Кольцо полиномов}

Пусть $K$~"--- кольцо, а $K[x]$~"--- кольцо полиномов. Напомню, что $K[x]$ состоит из последовательностей
\[
	(a_0, a_1, a_2, \dots),
\]
в которых только конечное число элементов $a_\nu \in K$ отлично от нуля. Такое множество обозначается символом $\bigoplus\limits_{\nu \geq 0} K$. Для удобства такие последовательности записываются в виде
\[
	a_0 x^0 + a_1 x^1 + a_2 x^2 + \dots,
\]
где только конечное число слагаемых отлично от нуля. Переменная~$x$ считается фиктивной, в неё априори ничего не <<подставляется>>, за исключением случаев, когда подстановка вместо $x$ числа $q \in K$ является \emph{алгебраической} операцией~"--- то есть вычисляется за \emph{конечное} число сложений и умножений. (В случае с полиномами это всегда так, а в случае с рядами~"--- нет.)

\begin{definition}
	\emph{Старшей степенью} полинома $f \in K[x]$ называется наибольшая степень $x^\nu$, для которой $a_\nu \neq 0$.
	% \label{<+label+>}
\end{definition}

\begin{definition}
	\emph{Старшим коэффициентом} полинома $f \in K[x]$ называется коэффициент $a_\nu$ при старшей степени.
	% \label{<+label+>}
\end{definition}

\begin{definition}
	\emph{Степенью} $\deg{f}$ полинома $f \in K[x]$ называется показатель старшей степени.
	% \label{<+label+>}
\end{definition}

Операция сложения определена покомпонентно:
\[
	+\colon
	K[x] \times K[x]
	\to
	K[x]
\]
действует так:
\[
	(a_0, a_1, a_2, \dots)
	+
	(b_0, b_1, b_2, \dots)
	:=
	(a_0 + b_0, a_1 + b_1, a_2 + b_2, \dots)
	.
\]
В обозначениях с переменной~$x$
\[
	\left( 
		\sum_{i = 0}^{m}
			a_i x^i
	\right)
	+
	\left( 
		\sum_{j = 0}^{n}
			a_j x^j
	\right)
	=
	\sum_{k = 0}^{\max(m, n)}
		(a_k + b_k) x^k
\]
(мы полагаем $a_\mu = 0$ для всех $\mu > m$ и $b_\nu = 0$ для всех $\nu > n$).

Операция умножения на число~$k \in K$~"--- тоже покомпонентно:
\[
	\cdot \colon
	K \times K[x]
	\to
	K[x]
\]
действует так:
\[
	k \cdot
	(a_0, a_1, a_2, \cdot)
	:=
	(ka_0, ka_1, ka_2, \cdot)
	.
\]
В полиномиальных обозначениях:
\[
	k \cdot
	\left( 
		\sum_{i = 0}^m
			a_i x^i
	\right)
	:=
	\sum_{i = 0}^m
		ka_i x^i
	.
\]

А вот умножение полиномов
\[
	\cdot \colon
	K[x] \times K[x] \to K[x]
\]
--- это та операция, которую легче записать в полиномиальных обозначениях (раскрываем скобки по дистрибутивности):
\[
	\left( 
		\sum_{i = 0}^{m}
			a_i x^i
	\right)
	\left( 
		\sum_{j = 0}^{n}
			b_j x^j
	\right)
	:=
	\sum_{i = 0}^m
	\sum_{j = 0}^n
		a_i b_j x^{i+j}
	.
\]
Запишем этот полином, указав явно коэффициент при~$x^s$. Для этого положим $s := i+j$. Суммирование по целочисленному прямоугольнику $([0, m] \cap \NN) \times ([0, n] \cap \NN)$ тогда перейдёт в суммирование по диагональным прямым, задаваемым уравнением $s = i + j$ при фиксированном~$s$ (нарисуй!). Полагая, как водится, $a_\mu = 0$ при $\mu > m$ и $\mu < 0$ и $b_\nu = 0$ при $\nu > n$ и $\nu < 0$, можно записать
\[
	\sum_{i = 0}^m
	\sum_{j = 0}^n
		a_i b_j x^{i+j}
	=
	\sum_{s = 0}^{m+n}
		\left( 
			\sum_{i + j = s}
				a_i b_j
		\right)
		x^s
		.
\]
Это и есть искомая формула. В обозначениях с последовательностями получаем
\[
	(a_0, a_1, a_2, \dots)
	\cdot
	(b_0, b_1, b_2, \dots)
	:=
	( 
		a_0 b_0
		,\ 
		a_0 b_1 + a_1 b_0
		,\ 
		a_0 b_2 + a_1 b_1 + a_2 b_0
		,\ 
		\dots
	).
\]
Результирующее выражение называется \emph{свёрткой} последовательностей~$(a_\mu)$ и~$(b_\nu)$.

\subsection{Гомоморфизм вычисления}

В полином $f(x) \in K[x]$ можно подставить число $a \in K$ вместо~$x$. Получившееся в результате число $f(a) \in K$ называется \emph{значением} полинома~$f$ на элементе~$a$.

Если есть какое-то явление, полезно посмотреть на него с функиональной точки зрения, то есть определить функцию, которая это делает, и посмотреть на её свойства.

Определим \emph{гомоморфизм вычисления на элементе~$a \in K$}:
\[
	\ev_a \colon K[x] \to K
	\colon f(x) \mapsto f(a)
	.
\]

\begin{lemma}[деление полиномов с остатком]
	Для любых многочленов~$f,\ g \in \kk[x]$ с коэффициентами в \emph{поле}~$\kk$ существует \emph{единственная} пара многочленов $q,\ r \in \kk[x]$, такая, что
	\[
		f = g \cdot q + r
	\]
	и либо $\deg{r} < \deg{g}$, либо $r = 0$.
	% \label{<+label+>}
\end{lemma}

\begin{lemma}[вычисление значения полинома в точке]
	Остаток от деления многочлена
	\[
		f(x) = a_0 + a_1 x + \dots + a_n x^n
	\]
	на двучлен $x - \alpha$~"--- это константа, равная значению~$f(\alpha)$ многочлена~$f$ при $x = \alpha$.
	% \label{<+label+>}
\end{lemma}
\begin{proofm}
	Поделим $f(x)$ на $x - \alpha$ с остатком:
	\[
		f(x) = (x - \alpha) \cdot g(x) + r(x)
		.
	\]
	Так как $\deg{r} < \deg(x - \alpha) = 1$, то $\deg{r} = 0$, и $r$ на самом деле константа.

	Теперь применим к обеим частям гомоморфизм~$\ev_{\alpha}$:
	\begin{align*}
		f(\alpha) &= (\alpha - \alpha) \cdot g(\alpha) + r,
		\\
		f(\alpha) &= r.
		\proofmark
	\end{align*}
\end{proofm}
\begin{remark}
	Вычисление остатка от деления $f(x)$ на $(x - \alpha)$~"--- значительно более быстрый способ вычисления~$f(\alpha)$, чем подстановка напрямую.
	% \label{<+label+>}
\end{remark}

\begin{definition}
	Элемент~$\alpha \in K$ называется \emph{корнем} многочлена~$f \in K[x]$, если $f(\alpha) = 0$.
	% \label{<+label+>}
\end{definition}
\begin{remark}
	Это равносильно тому, что остаток от деления $f(x)$ на $(x - \alpha)$ равен нулю. А это значит, что $f(x)$ делится на $(x - \alpha)$.
	% \label{<+label+>}
\end{remark}

\begin{lemma}
	Пусть $K$~"--- целостное кольцо, а $f$ имеет различные корни $(\alpha_1, \dots, \alpha_p)$. Тогда $f$ делится на произведение $\prod\limits_{i = 1}^{p} (x - \alpha_i)$.
	% \label{<+label+>}
\end{lemma}
\begin{proofm}
	Так как в~$K$ нет делителей нуля и $\alpha_i - \alpha_j \neq 0$ при $i \neq j$, применим предыдущую лемму: $\alpha_1$~"--- корень~$f$, так что существует~$q(x) \in K[x]$, такой, что
	\[
		f(x) = q(x) \cdot (x - \alpha_1).
	\]
	Применяя $\ev_{\alpha_i}$ к этому равенству, получаем, что $(\alpha_2, \dots, \alpha_p)$~"--- корни~$q(x)$. Остаётся применить индукцию и получить, что
	\[
		f(x) = \prod_{i = 1}^{p} (x - \alpha_p).
		\proofmark
	\]
\end{proofm}

\begin{corrolary}
	Если у $f$ есть $p$ различных корней, то $\deg{f} \geq p$.
\end{corrolary}

\begin{definition}
	Пусть $\kk$~"--- поле. Число $\alpha \in \kk$ называется \emph{$m$-кратным корнем многочлена~$f \in \kk[x]$}, если существует многочлен~$q \in \kk[x]$, такой, что
	\[
		f(x) = 
		(x - \alpha)^m \cdot q(x)
	\]
	и $q(\alpha) \neq 0$.

	Корни кратности~$m \geq 2$ называются \tdef{кратными}.
	% \label{<+label+>}
\end{definition}

\begin{remark}
	Так как в кольце полиномов априори не задана топология (но её можно задать; этим мы заниматься не будем) или метрика, то нельзя говорить о производной многочлена как оператора в смысле анализа. Однако есть способ ввести оператор производной~$D \colon K[x] \to K[x]$, исходя из чисто алгебраических соображений, так, чтобы выполнялись привычные из анализа свойства: для любых $f,\ g \in K[x]$ и любого $\alpha \in K$
	\begin{align*}
		D(f + g) &= D(f) + D(g),
		\\
		D(\alpha f) &= \alpha \, D(f),
		\\
		D(f \cdot g) &= D(f) \cdot g + f \cdot D(g),
		\\
		D(x^m) &= mx^{m-1}
		\rlap{$
			\quad
			\text{при $m > 0$,}
		$}
		\\
		D(\alpha) &= 0.
	\end{align*}
	% \label{<+label+>}
\end{remark}

\begin{lemma}
	Число $\alpha \in K$ является кратным корнем полинома~$f \in K[x]$ тогда и только тогда, когда
	\[
		f(\alpha) =
		f'(\alpha) =
		0
		.
	\]
	\label{<+label+>}
\end{lemma}
\begin{proofm}
	Если $\alpha$~"--- кратный корень~$f$, то
	\[
		f(x) = (x - \alpha)^2\,g(x),
	\]
	где, возможно, $g(\alpha) = 0$. Продифференцируем это выражение и подставим~$x = \alpha$:
	\begin{align*}
		f'(x) &= 2(x - \alpha)\,g(x) + (x - \alpha)^2\,g'(x)
		\\
		f'(\alpha) &= 2(\alpha - \alpha)\,g(\alpha) + (\alpha - \alpha)^2\,g'(\alpha) = 0
		.
	\end{align*}
	Если же~$\alpha$~"--- не кратный корень, то
	\[
		f(x) = (x - \alpha)\,h(x),
	\]
	где $h(\alpha) \neq 0$. Тогда
	\begin{align*}
		f'(x) &= h(x) + (x - \alpha)\,h'(x)
		\\
		f'(\alpha) &= h(\alpha) \neq 0.
		\proofmark
	\end{align*}
\end{proofm}

\begin{definition}
	Многочлен $f \in \kk[x]$ называется \emph{сепарабельным}, если $\gcd{f, f'} = 1$. Это означает, что $f$ не имеет кратных корней ни в каком кольце $K \supset \kk$.
	\label{<+label+>}
\end{definition}

\begin{lemma}
	Пусть $\kk$~"--- произвольное поле. Для любого набора многочленов $f_1, \dots, f_n \in \kk[x]$ существует единственный приведённый многочлен~$d \in \kk[x]$, который делит каждый из многочленов~$f_\nu$ и делится на любой многочлен, делящий каждый из~$f_\nu$. При этом~$d$ представляется в виде
	\[
		d = h_1 f_1 + \dots + h_n f_n
	\]
	для некоторых $h_\nu \in \kk[x]$, то есть в виде $\kk[x]$-линейной комбинации векторов $f_1, \dots, f_n$. Иными словами, $d \in (f_1, \dots, f_n)$, где
	\[
		(f_1, \dots, f_n) :=
		\set{
			h_1 f_1 + \dots + h_n f_n
		}{h_\nu \in \kk[x]}
	\]
	--- идеал, порождённый над $\kk[x]$ полиномами $f_1$, \dots, $f_n$.
	
	При этом произвольно взятый многочлен~$g \in \kk[x]$ представ\'{и}м в виде линейной комбинации полиномов из набора~$f_1$, \dots, $f_n$ тогда и только тогда, когда он делится на~$d$.
	\label{<+label+>}
\end{lemma}

Такой многочлен называется \emph{наибольшим общим делителем} многочленов $f_1, \dots, f_n$ и обозначается~$\gcd{f_1, \dots, f_n}$.



\begin{remark}
	Сравни это определение с определением наибольшего общего делителя набора чисел~$n_1, \dots, n_k$ как наименьшего положительного числа, находящегося в идеале
	\[
		(n_1, \dots, n_k) :=
		\set{a_1 n_1 + \dots + a_k n_k}{a_i \in \ZZ}
		.
	\]
	\label{<+label+>}
\end{remark}

\begin{definition}
	Полином $f \in K[x]$ называется \emph{приведённым}, если коэффициент при старшей степени равен~$1$: $a_{\deg{f}} = 1$.
\end{definition}

\begin{definition}
	Многочлен~$f \in K[x]$ с коэффициентами в целостном кольце~$K$ называется \emph{неприводимым}, если из $f = gh$ следует, что $g$ или $h$~"--- обратимая константа.
	\label{<+label+>}
\end{definition}

\subsection{Факторкольцо полиномов}

Пусть $\kk$~"--- поле. Факторкольцо $\kk[x]/(f)$, где $f \in \kk[x]$, устроено аналогично кольцу $\ZZ/(n)$.

Пусть $f \in \kk[x]$~"--- произвольный фиксированный полином, отличный от константы. Обозначим
\[
	(f) := \set{hf}{h \in \kk[x]}
\]
идеал всех многочленов, делящихся на~$f$. Отношение
\[
	\bigl(
		g_1 \sim g_2
	\bigr)
	:=
	\bigl(
		g_1 - g_2 \in (f)
	\bigr)
\]
является отношением эквивалентности (почему?) и разбивает~$\kk[x]$ в объединение непересекающихся классов
\[
	[g] :=
	g + (f) =
	\set{g + hf}{h \in \kk[x]},
\]
которые называются \emph{классами вычетов} по модулю~$f$. Сложение и умножение классов вычетов задаётся формулами
\begin{align*}
	[g] + [h] &:= [g + h]
	\\
	[g] \cdot [h] &:= [g \cdot h]
\end{align*}

Нуль кольца $\kk[x]/(f)$~"--- это класс $[0] = (f)$; единица~"--- это класс $[1] = 1 + (f)$. Так как константы не делятся на многочлены положительной степени, то классы всех констант из~$\kk$ различны по модулю~$f$, и поэтому для $c \in \kk$ вместо $[c]$ можно писать~$c$.

Так как любой полином $g \in \kk[x]$ единственным образом раскладывается в виде
\[
	g = hf + r,
	\rlap{%
		\qquad
		где
		$\deg{r} < \deg{f}$,
	}
\]
то в каждом классе~$[g]$ имеется единственный представитель $r \in [g]$, такой, что $\deg{r} < \deg{f}$.

Иными словами, каждый класс $[g] \in \kk[x]/(f)$ однозначно записывается в виде
\begin{align*}
	[g(x)] &=
	[a_0 + a_1 x + \dots + a_{n-1} x^{n-1}]
	\nlinea{=}
	a_0 + a_1 [x] + \dots + a_{n-1} [x]^{n-1}
	\nlinea{=}
	a_0 + a_1 \theta + \dots + a_{n-1} \theta^{n-1}
	,
\end{align*}
где $a_i \in \kk$ и $\theta := [x]$.

При этом класс~$\theta = [x]$ удовлетворяет в кольце~$\kk[x]/(f)$ уравнению $f(\theta) = 0$, так как
\[
	f(\theta)
	=
	f\bigl( [x] \bigr)
	=
	[f(x)]
	=
	[0]
	.
\]

Поэтому сложение и умножение классов в~$\kk[x]/(f)$ можно формально интерпретироваьт как сложение и умножение <<многочленов>> вида
\[
	a_0 + a_1 \theta + \dots + a_{n-1} \theta^{n-1}
\]
по стандартным правилам раскрытия скобок и приведения подобных с учётом того, что~$\theta$ удовлетворяет соотношению $f(\theta) = 0$. Поэтому часто обозначают
\[
	\kk[\theta] := \frac{\kk[x]}{(f)},
\]
где $f(\theta) = 0$, и называют \tdef{расширением} поля~$\kk$ посредством \emph{присоединения} к нему корня~$\theta$ многочлена~$f \in \kk[x]$.

\begin{example}
	Кольцо $\QQ[x]/(x^2 - 2)$ можно воспринимать как множество формальных записей вида $a + b\theta$, где $\theta := [x]$. Сложение и умножение таких записей происходит по обычным правилам с учётом того, что \mbox{$\theta^2 = 2$}. По этой причине можно \emph{определить} $\sqrt{2} := \theta := [x]$ и воспринимать факторкольцо $\QQ[x]/(x^2 - 2)$ как множество записей вида $a + b\sqrt{2}$.

	Например, перемножим $a + b\sqrt{2}$ и $c + d\sqrt{2}$:
	\[
		\bigl(
			a + b\sqrt{2}
		\bigr)
		\bigl(
			c + d\sqrt{2}
		\bigr)
		=
		ac + ab \sqrt{2} + bc \sqrt{2} + bd \bigl( \sqrt{2} \bigr)^2
		=
		ac + 2bd + (ab + bc)\,\sqrt{2},
	\]
	потому что $\bigl( \sqrt{2} \bigr)^2 := 2$ (обрати внимание, что~$\sqrt{2}$ в данном случае~"--- \emph{формальная переменная с определёнными свойствами}, а не число).
\end{example}

\begin{lemma}
	Пусть $\kk$~"--- произвольное поле. Кольцо $\kk[x]/(f)$ является полем тогда и только тогда, когда $f$ неприводим в~$\kk[x]$.
	\label{<+label+>}
\end{lemma}
\begin{proof}
	Если $f = gh$, где степени $g$ и $h$ строго меньше степени~$f$, то классы $[g] \neq [0]$ и $[h] \neq [0]$ будут делителями нуля в~$\kk[x]/(f)$. Но в поле нет делителей нуля.

	Если $f$ неприводим, то для любого $g \notin (f)$ имеем $\gcd{f, g} = 1$, так что $fh + gq = 1$ для некоторых $h$ и~$q \in \kk[x]$, откуда
	\begin{align*}
		[fh + gq] &= [1],
		\\
		[f][h] + [g][q] &= [1],
		\\
		[0] + [g][q] &= [1],
		\\
		[g][q] &= [1],
	\end{align*}
	и произвольный
	$
		[g] \in
		\bigl( \kk[x]/(f) \bigr)
		^\times
		% \setminus \set*{[0]}
	$ обратим.
\end{proof}

% \begin{example}
	Напишем явную формулу для вычисления обратного элемента к числу $a := a_0 + a_1 \theta$ в поле~$\QQ[\theta]$, где $\theta^2 + \theta + 1 = 0$.

	Это равносильно нахождению обратного элемента в поле
	\[
		\frac{\QQ[x]}{(x^2 + x + 1)}
		.
	\]

	Пусть $b := b_0 + b_1\theta$ --- обратный к~$a$. Тогда $ab = 1$, то есть
	\begin{align*}
		(a_0 + a_1 \theta) (b_0 + b_1 \theta) &= 1
		\\
		a_0 b_0 + (a_1 b_0 + a_0 b_1) \theta + a_1 b_1 \theta^2 &= 1
		.
	\end{align*}
	Так как $\theta^2 + \theta + 1 = 0$, то $\theta^2 = -\theta - 1$ (стершую степень многочлена, аннулирующего~$\theta$, можно выразить через его меньшие степени). Тогда
	\begin{align*}
		a_0 b_0 + (a_1 b_0 + a_0 b_1) \theta + a_1 b_1 (-\theta - 1) &= 1
		\\
		(a_0 b_0 - a_1 b_1) +
		(a_1 b_0 + a_0 b_1 - a_1 b_1) \theta &= 1
	\end{align*}
	Приравняем коэффициенты при одинаковых степенях~$\theta$:
	\[
		\begin{system}
			a_0 b_0 - a_1 b_1 &= 1,
			\\
			a_1 b_0 + (a_0 - a_1) b_1 &= 0.
		\end{system}
	\]
	В матричной записи
	\[
		\begin{pmatrix}
			a_0 & -a_1
			\\
			a_1 & a_0 - a_1
		\end{pmatrix}
		\Vect{b_0 \\ b_1}
		=
		\Vect{1 \\ 0}.
	\]
	Так как мы находимся в поле, то всякий ненулевой элемент обратим. Решим систему методом Крамера:
	\[
		\Delta :=
		\begin{determ}
			a_0 & -a_1
			\\
			a_1 & a_0 - a_1
		\end{determ}
		=
		a_0 (a_0 - a_1) + a_1^2 =
		a_0^2 - a_0 a_1 + a_1^2 =
		(a_0 - a_1)^2 + a_0 a_1.
	\]

	\textstars

	Есть другой способ.

	Достаточно найти обратные к элементам вида~$\theta - a$. Это можно сделать по алгоритму Евклида: если $h(\theta)$ --- класс, обратный к $\theta - a$, то $h(\theta)$ задаётся таким многочленом $h \in \QQ[x]$, что $[h(x)][x-a] = [1]$, то есть $[h(x)(x-a)] = [1]$, то есть $h(x)(x-a) + \bigl( (x^2 + x + 1) \bigr) = 1 + \bigl( (x^2 + x + 1) \bigr)$, то есть
	\[
		h(x) (x - a) + g(x)(x^2 + x + 1) = 1
	\]
	для некоторого $g \in \QQ[x]$.

	Такое разложение можно найти по алгоритму Евклида. Пусть $E_0 := x^2 + x + 1$, $E_1 := x-a$. Найдём $E_2 := E_0 \bmod E_1$ (поделим столбиком):
	\[
		x^2 + x + 1 =
		(x-a) \cdot x +
		\underbrace{%
			\bigl( (a+1)x + 1 \bigr)
		}_{E_2}
	\]
	Итак, $E_0 = xE_1 + E_2$, откуда $E_2 = E_0 - xE_1$; далее найдём $E_3 := E_1 \bmod E_2$:
	\[
		x-a = 
		\bigl( (a+1)x + 1 \bigr)
		\cdot
		\frac{1}{a+1}
		+
		\underbrace{%
			\bigl(
				-a - \frac{1}{a+1}
			\bigr)
		}_{E_3}
	\]
	Итак,
	\begin{align*}
		E_3 &=
		E_1 - \frac{E_2}{a + 1}
		\nlinea{=}
		E_1 - \frac{E_0 - xE_1}{a + 1}
		\nlinea{=}
		\frac{-E_0 + (a+1+x)E_1}{a + 1}
	\end{align*}
	Или, что то же самое,
	\begin{align*}
		-a - \frac{1}{a+1} &=
		\frac{-E_0 + (a+1+x)E_1}{a + 1}
		\\
		-a(a+1) - 1 &=
		-E_0 + (a+1+x)E_1
		\\
		1 &=
		\frac{-E_0}{-a^2 - a - 1} +
		\frac{(a+1+x)E_1}{-a^2 - a - 1}
		\\
		1 &=
		\frac{-E_0}{-a^2 - a - 1} +
		\frac{(a+1+x)E_1}{-a^2 - a - 1}
	\end{align*}

\end{example}



\begin{theorem}
	Для любого поля~$\kk$ и любого многочлена~$f \in \kk[x]$ существует такое поле~$\FF \supset \kk$, что $f$ раскладывается в $\FF[x]$ в произведение~$\deg{f}$ линейных множителей.
	\label{<+label+>}
\end{theorem}
\begin{proof}
	Докажем по индукции по степени $n = \deg{f}$.
	\begin{enumerate}
		\item Для многочлена~$f(x) = x - a$ степени $n = 1$ мы уже умеем строить такое поле~"--- это само поле~$\kk$.

		\item Теперь пусть $n \in \NN$ и для всех многочленов степени, меньшей $n$, мы умеем строить нужное поле.
			\begin{itemize}
				\item Если $\deg{f} = n$ и $f$ приводим, то $f = gh$, где $\deg{g} < n$, $\deg{h} < n$. Построим поле $\LL \supset \kk$, над которым $g$ полностью раскладывается на линейные множители, а затем поле $\FF \supset \LL$, над которым разложится и~$h$.

				\item Tеперь пусть $f$ неприводим. Рассмотрим поле $\LL := \kk[x]/(f)$. Оно содержит~$\kk$ в качестве классов констант, и~$f$ делится в~$\LL[x]$ на~$x - \theta$, где $\theta = x \bmod f$: $f(x) = (x - \theta) q(x)$ в~$\LL[x]$, причём $\deg{q} = n - 1$, и по индукции~$q$ раскладывается на линейные множители над некоторым полем~$\FF \supset \LL$.
			\end{itemize}
	\end{enumerate}
	Теорема доказана.
\end{proof}

\section{Расширения полей}

Пусть $\KK$ и~$\FF$ --- поля.

\begin{definition}
	Поле~$\KK$ называется \tdef{расширением} поля~$\FF$, если $\KK \supset \FF$.
\end{definition}

\begin{proposition}
	Если~$\KK \supset \FF$, то $\KK$ можно представлять себе как векторное пространство над~$\FF$.
\end{proposition}

\begin{definition}
	Если~$\KK \supset \FF$, то размерность~$\dim_{\FF}{\KK}$ поля~$\KK$ как векторного пространства над~$\FF$ называется \tdef{степенью} расширения и обозначается~$[\KK : \FF]$:
	\[
		[\KK : \FF] := \dim_{\FF}{\KK}
		.
	\]
\end{definition}

\begin{example}
	Поле $\CC := \RR[t]/(t^2 + 1)$ --- это векторное поле над~$\RR$ размерности~$\dim_{\RR}{\CC} = [\CC : \RR] = 2$.
\end{example}

\begin{definition}
	Последовательность
	\[
		\KK \supset \LL \supset \FF
	\]
	расширений поля~$\FF$ называется \tdef{башней расширений}.
\end{definition}

\begin{lemma}
	Если $\KK \supset \LL \supset \FF$ --- башня расширений, то
	\[
		[\KK : \LL]
		\cdot
		[\LL : \FF]
		=
		[\KK : \FF]
		.
	\]
\end{lemma}
\begin{proof}
	Требуется доказать среднее равенство в цепочке
	\[
		[\KK : \LL] \cdot [\LL : \FF] =
		\dim_{\LL}{\KK} \cdot \dim_{\FF}{\LL} =
		\dim_{\FF}{\KK}
		=
		[\KK : \FF]
		.
	\]
	Чтобы рассуждать о размерности, нужен базис. Пусть
	\begin{align*}
		\bu &:= (u_1, \dots, u_m) \subset \KK,
		\\
		\bv &:= (v_1, \dots, v_n) \subset \LL
	\end{align*}
	--- базисы $\KK$ как векторного пространства над~$\LL$ и~$\LL$ как векторного пространства над~$\FF$. Тогда
	\[
		\bu \otimes \bv :=
		\set{u_i \otimes v_j}{%
			\begin{gathered}
				1 \leq i \leq m,
				\\
				1 \leq j \leq n
			\end{gathered}
		}
		=
		\set{u_i v_j}{%
			\begin{gathered}
				1 \leq i \leq m,
				\\
				1 \leq j \leq n
			\end{gathered}
		}
	\]
	--- базис $\KK$ над~$\FF$, так как если
	\[
		k = a_1 u_1 + \dots + a_m u_m =
		\sum_{i = 1}^m
			a_i u_i
		,
	\]
	где $a_i \in \LL$ --- разложение~$k \in \KK$ над~$\LL$, а элемент~$a_i \in \LL$ раскладывается над~$\FF$ в виде
	\[
		a_i =
		b_{i1} v_1 + \dots + b_{in} v_n =
		\sum_{j = 1}^n
			b_{ij} v_j
		,
	\]
	где $b_{ij} \in \FF$,
	% а верхние индексы не обозначают степени,
	то
	\[
		k =
		\sum_{i = 1}^m
			a_i u_i
		=
		\sum_{i = 1}^m
		\left( 
			\sum_{j = 1}^n
				b_{ij} v_j
		\right)
		u_i
		=
		\sum_{i = 1}^m
		\sum_{j = 1}^n
				b_{ij} \cdot v_j u_i
		.
	\]
	Итак, каждый элемент~$k \in \KK$ раскладывается в~$\FF$ по базису~$\bu \otimes \bv$, причём в этом базисе
	\[
		mn =
		\deg_{\LL}{\KK} \cdot \deg_{\FF}{\LL}
		=
		[\KK : \LL] \cdot [\LL : \FF]
	\]
	элементов. С другой стороны, $\dim_{\FF}{\KK} = \dim(\bu \otimes \bv) = mn$.
\end{proof}

\subsection{Конечные поля}

\begin{definition}
	\tdef{Конечное поле} --- это поле, состоящее из конечного числа элементов.
	% \label{<+label+>}
\end{definition}

\begin{definition}
	\[
		\abs{\FF} :=
		\left( 
			\parbox{1.2in}{%
				\centering
				число элементов в~$\FF$%
			}
		\right)
		.
	\]
	% \label{<+label+>}
\end{definition}

\begin{example}
	Кольцо~$\ZZ/(p)$ при простом~$p$ является полем.
\end{example}

\begin{remark}
	Мы уже видели, что при составном~$n \in \ZZ$ в кольце~$\ZZ/(n)$ есть делители нуля, и поэтому оно не является полем.
	% \label{<+label+>}
\end{remark}

\begin{lemma}
	Если поле~$\FF$ конечно, то~$\abs{\FF} = p^n$, где~$p$ --- простое, а~$n \in \NN$.
	% \label{<+label+>}
\end{lemma}
\begin{proof}
	% Пусть $\abs{\FF} = m \in \NN$. Так как~$\FF^\times = \FF \setminus \set*{0}$ образует конечную группу по умножению, то в ней есть образующая~$a \in \FF$, то есть такой элемент, что $a^m = 1$ и $a^\mu \neq 1$ при~$\mu < m$. 

	% Если $m = nk$, где $n$, $k \in \NN$, то
	% \[
	% 	
	% \]
\end{proof}

\begin{remark}
	Хоть кольцо~$\ZZ/(p^m)$ состоит из~$\abs{\ZZ/(p^m)} = p^m$ элементов, но при~$m \neq 1$ оно не является полем, так что конечные поля так не построить.

	Вместо этого стартуем с поля~$\FFp$ и строим его расширения, используя процедуру присоединения корней:
	\begin{align*}
		\FFp[2] &:= \frac{\FFp{}[x]}{(f_2)} = \FFp{}[\theta_2],
		\\
		\FFp[3] &:= \frac{\FFp{}[x]}{(f_3)} = \FFp{}[\theta_3],
		\\
		\hbox to 3em
		{\dotfill}
		&
		\mkern -6.2mu
		\hbox to 2in
		{\dotfill}
		\\
		\FFp[n] &:= \frac{\FFp{}[x]}{(f_n)} = \FFp{}[\theta_n],
	\end{align*}
	где~$f_k \in \FFp{}[x]$ --- неприводимый полином степени~$k$, а~$\theta_k$ удовлетворяет уравнению~$f_k(\theta_k) = 0$.
	
	% \label{<+label+>}
\end{remark}

\subsubsection{Гомоморфизм Фробениуса}

\begin{definition}
	Гомоморфизм Фробениуса --- это гомоморфизм полей
	\[
		\Phi \colon
		\FF_q \to \FF_q
		\colon
		x \mapsto x^{q}
		.
	\]
\end{definition}

% \begin{lemma}
% 	В поле характеристики~$p$
% 	\[
% 		(a + b)^p = a^p + b^p
% 		.
% 	\]
% \end{lemma}
% \begin{proof}
% 	Раскроем скобки:
% 	\[
% 		(a + b)^p =
% 		\sum_{k = 0}^p
% 			\binom{p}{k}
% 			a^k b^{p-k}
% 			.
% 	\]
% 	Так как
% 	\[
% 		\binom{p}{k} =
% 		\frac{p!}{k!\,(p-k)!}
% 		=
% 		p \cdot w
% 		,
% 	\]
% 	где $w = (p-1)!/\bigl( k!\,(p-k)! \bigr)$, то $p$ делит~$\binom{p}{k}$, и поэтому~$\binom{p}{k} = 0$ (все числа, кратные~$p$, в поле )<++>
% \end{proof}<++>





% \subsection{Поля}
% 
% \begin{definition}
% 	Кольцо~$F$ называется \emph{полем}, если всякий ненулевой элемент в нём обратим.
% 	\label{<+label+>}
% \end{definition}<++>
